
\textbf{Corollaire 5.4.} \uline{La filtration induite par} $K^{\bullet}(X)$ \uline{sur} $K^{\bullet}(S)$ \uline{est identique à la filtration de} $K^{\bullet}(S)$.

\textbf{Corollaire 5.5.} \uline{Soient} $S$ \uline{un schéma quasi}-\uline{compact}, $\underline{E}_1,\ldots,\underline{E}_n$ \uline{des} $\underline{O}_S$-\uline{Modules} \uline{localement libres de type fini}, $\mathbf D(\underline{E}_i)$ \uline{le schéma de drapeaux de} $\underline{E}_i$, $X=\mathbf D(\underline{E}_1)\times_S\ldots\times_S \mathbf D(\underline{E}_n)$. \uline{Alors l'homomorphisme canonique} :
\[
  \Gr K^{\bullet}(S)\longrightarrow \Gr K^{\bullet}(X)
\]
\uline{est injectif}.

Le raisonnement de 3.1 permet de se ramener à montrer la même propriété quand $X$ est un fibré projectif sur $S$. Ce cas résulte alors immédiatement de 5.4.

\textbf{Corollaire 5.6.} \uline{Avec les notations et les hypothèses de} 5.1, $\Gr K^{\bullet}(X)$ \uline{est un} $\Gr K^{\bullet}(S)$-\uline{module libre}, \uline{admettant pour base} $1,\bar x,\ldots,\bar x^r$. \uline{Pour tout} $k$, \uline{on a} :
\[
  \Gr^kK^{\bullet}(X)=\Gr^kK^{\bullet}(S)\oplus \Gr^{k-1}K^{\bullet}(S).\bar x\oplus\ldots\oplus \Gr^{k-r}K^{\bullet}(S).\bar x^r .
\]
\uline{Enfin}, \uline{les éléments} $1,\bar x,\ldots,\bar x^{r+1}$ \uline{sont liés par la relation} :
\[
  \sum_{i=0}^{r+1}c^{r+1-i}(\check E)\bar x^i=0,
\]
\uline{où} $c^k(\check E)$ \uline{est la} $k$-\uline{ième classe de Chern de} $\check E$ (V 6.7).

Conséquence immédiate de 1.1, 5.3, et (1.13.1).

\textbf{Corollaire 5.7.} \uline{Avec les notations et les hypothèses de} 4.5, \uline{la} $\Gr K^{\bullet}(S)$-\uline{algèbre} $\Gr K^{\bullet}(X)$ \uline{est engendrée par les} $\bar c_i^{(j)}$, \uline{où} $\bar c_i^{(j)}$ \uline{est la classe dans} $\Gr^iK^{\bullet}(X)$ \uline{de} $\gamma^i(F_j-p_j)$. \uline{L'idéal des relations entre les} $\bar c_i^{(j)}$ \uline{est engendré par les relations déduites de} :
\[
  c(E-n)=\prod_{j=1}^{k}c(F_j-p_j),
\]
